\chapter{Le marché des options de taux}\label{ch:m2:the-rates-options-market}

Le 11 juin 2014, la Banque centrale européenne abaisse le taux qu'elle verse sur les dépôts
à $-0.10\%$, puis à $-0.50\%$ au cours des cinq années suivantes. Les rendements allemands à dix ans
sont restés négatifs pendant 38 mois entre le milieu de 2016 et le début de 2022. Pour le marché
des options, ce n'était pas une curiosité mais une rupture. Le modèle longtemps standard pour
ces options, celui de Black, suppose qu'un
taux évolue proportionnellement à son niveau, et un taux de $0.10\%$ qui passe à
$0.20\%$ a progressé de 100~\%~; sous zéro, le modèle ne fournit plus aucune réponse. Le marché
s'est tourné vers une cotation dans l'unité où raisonne un trader de taux, les points
de base de variation, avec le \omterm{def:m2:the-rates-options-market:normal}{modèle de Bachelier}, écrit en 1900 par un doctorant français.
Ce chapitre présente les options négociées sur les taux d'intérêt, les \omterm{def:m2:the-rates-options-market:cap}{caps}, les \omterm{def:m2:the-rates-options-market:cap}{floors} et
les \omterm{def:m2:the-rates-options-market:swaption}{swaptions}, la \omterm{def:m2:the-rates-options-market:normal}{volatilité normale} dans laquelle elles sont cotées, et les flux
structurels, issus des \omterm{def:m2:the-rates-options-market:swaption}{obligations remboursables par anticipation} et des crédits immobiliers, qui fixent leurs prix.

\section{Caps et floors}

\begin{definition}[Caplet, cap et floor]\label{def:m2:the-rates-options-market:cap}
Un \emph{caplet}\index{caplet} de prix d'exercice $K$ sur un taux variable pour une période de
durée $\delta$ verse $N\delta\max(L - K, 0)$ à la fin de la période, où $L$
est le taux fixé pour cette période. Un \emph{cap}\index{cap} est une succession de
caplets sur des périodes consécutives~; un \emph{floor}\index{floor} est la même succession de
floorlets, chacun versant $N\delta\max(K - L, 0)$.
\end{definition}

Un emprunteur qui paie un taux variable achète un \omterm{def:m2:the-rates-options-market:cap}{cap} pour limiter ce qu'il peut payer~; un
prêteur, ou un investisseur en obligations à taux variable, achète un \omterm{def:m2:the-rates-options-market:cap}{floor}. Acheter un \omterm{def:m2:the-rates-options-market:cap}{cap} et
vendre un \omterm{def:m2:the-rates-options-market:cap}{floor} constitue un collar, souvent fixé de sorte que les deux primes se compensent. Chaque
\omterm{def:m2:the-rates-options-market:cap}{caplet} est une option sur un seul taux forward, si bien qu'un \omterm{def:m2:the-rates-options-market:cap}{cap} se valorise comme une somme de
options à une période, avec une volatilité par période, ou une seule volatilité
cotée pour l'ensemble du \omterm{def:m2:the-rates-options-market:cap}{cap}.

\begin{example}[Un cap et un collar à coût nul]\label{ex:m2:the-rates-options-market:cap}
Sur une courbe plate à 4~\%, un \omterm{def:m2:the-rates-options-market:cap}{cap} de cinq ans à 4.5~\% portant sur 100~millions de dollars
à taux annuel, la première période étant déjà fixée, comporte quatre \omterm{def:m2:the-rates-options-market:cap}{caplets} fixés entre un et
quatre ans. Avec une \omterm{def:m2:the-rates-options-market:normal}{volatilité normale} de 100 points de base par an, ils valent
182\,874, 310\,339, 401\,400 et 470\,709~dollars~: au total 1.37~million de dollars, soit 137
points de base du notionnel. Un \omterm{def:m2:the-rates-options-market:cap}{floor} à 3.5~\% vaut exactement autant, puisque le
taux forward se situe à mi-chemin des strikes et que le modèle normal est symétrique~:
l'emprunteur qui achète le \omterm{def:m2:the-rates-options-market:cap}{cap} et vend le \omterm{def:m2:the-rates-options-market:cap}{floor} ne paie rien et conserve son taux
entre 3.5~\% et 4.5~\%.
\end{example}

\section{Swaptions}

\begin{definition}[Swaption, payeuse, receveuse~; obligation remboursable par anticipation]\label{def:m2:the-rates-options-market:swaption}
Une \emph{swaption}\index{swaption} est une option permettant d'entrer dans un \omterm{def:m2:interest-rate-swaps:swap}{swap de taux}
à un taux fixe $K$ à une date future. Une \emph{swaption
payeuse}\index{swaption payeuse} confère le droit de payer le taux fixe~; une
\emph{swaption receveuse}\index{swaption receveuse} confère le droit de recevoir le taux fixe.
À l'échéance $T$, une swaption payeuse vaut $N A(T)\max(S_T - K, 0)$, où $S_T$ est le
\omterm{def:m2:government-bonds:zero}{taux au pair} du swap sous-jacent à cette date et $A(T)$ son annuité. Une \emph{obligation
remboursable par anticipation}\index{obligation remboursable par anticipation} (callable) est une obligation que son émetteur peut
rembourser avant l'échéance, à un prix et à des dates fixés.
\end{definition}

Les \omterm{def:m2:the-rates-options-market:swaption}{swaptions} se désignent par leur échéance et leur durée~: une 1y$\times$10y est une option à un an sur
un swap de dix ans. Une payeuse mise sur une hausse des taux, une receveuse sur une baisse~;
la combinaison fixée au taux forward, le straddle, parie sur l'ampleur des mouvements de taux,
dans un sens ou dans l'autre (\cref{fig:m2:the-rates-options-market:payoff}). Le
annuité convertit un écart de taux en argent~: la valorisation utilise le taux de swap forward
et l'annuité forward~; dans cette mesure, le taux de swap est une martingale, si bien
la \omterm{def:m2:the-rates-options-market:swaption}{swaption} se valorise comme une option sur un taux unique.

\begin{omfigure}
\begin{tikzpicture}
\begin{axis}[width=11cm,height=4.8cm,
  xlabel={taux swap à dix ans à l'échéance (\%)},ylabel={valeur pour 100 de notionnel},
  tick label style={font=\small},label style={font=\small},
  xmin=2,xmax=6,ymin=0,ymax=16,grid=major,grid style={gray!25},
  legend columns=3,legend style={font=\footnotesize,at={(0.5,-0.32)},anchor=north,draw=none,fill=none,/tikz/every even column/.append style={column sep=8pt}},
  table/col sep=comma]
\addplot[omThm,very thick] table[x=s,y=payer]{figdata/markets-2/13-the-rates-options-market/payoff.csv};
\addplot[omDef,very thick] table[x=s,y=receiver]{figdata/markets-2/13-the-rates-options-market/payoff.csv};
\addplot[omThm,thick,dashed] table[x=s,y=payer_value]{figdata/markets-2/13-the-rates-options-market/payoff.csv};
\legend{payeuse à l'échéance,receveuse à l'échéance,payeuse un an avant l'échéance}
\end{axis}
\end{tikzpicture}

\omcaption{\omterm{def:m2:the-rates-options-market:swaption}{Swaption payeuse} et \omterm{def:m2:the-rates-options-market:swaption}{swaption receveuse} 1y$\times$10y, fixées au taux forward de 4~\%,
pour 100 de notionnel, sur une courbe plate à 4~\%. À l'échéance, chacune verse l'annuité (7.80)
multipliée par l'écart du taux swap au-delà du strike~; un an avant, la payeuse
vaut en plus sa valeur temps, à une \omterm{def:m2:the-rates-options-market:normal}{volatilité normale} de 95 points de base.
À titre d'illustration~; données~: le tutoriel du chapitre.}\label{fig:m2:the-rates-options-market:payoff}
\end{omfigure}

Les \omterm{def:m2:the-rates-options-market:swaption}{obligations remboursables par anticipation} sont une source naturelle d'offre de \omterm{def:m2:the-rates-options-market:swaption}{swaptions}. Un émetteur qui
émet une obligation à dix ans remboursable par anticipation au bout de trois ans, détient le droit de cesser
de verser son coupon à partir de la troisième année, ce qui revient à une \omterm{def:m2:the-rates-options-market:swaption}{swaption receveuse}~: quand les taux baissent, il rembourse
l'obligation et se refinance. Les émetteurs qui se financent à taux variable swappent
l'obligation avec un dealer et cèdent ce droit au dealer dans le swap, que le
dealer peut annuler à la date de remboursement~; le dealer est alors long d'une
\omterm{def:m2:the-rates-options-market:swaption}{swaption receveuse}, donc longue de volatilité, et vend des \omterm{def:m2:the-rates-options-market:swaption}{swaptions} à d'autres investisseurs
pour transférer le risque (\cref{fig:m2:the-rates-options-market:callable}). Le
problème du week-end réalise cette opération.

\begin{omfigure}
\centering
\begin{tikzpicture}[font=\small,
  box/.style={draw,thick,rounded corners=2pt,align=center,minimum height=1cm,minimum width=2.4cm,inner sep=3pt},
  arr/.style={-{Stealth[length=2.5mm]},thick}]
\node[box,draw=omDef,fill=omDef!8] (inv) at (0,0) {investisseurs obligataires};
\node[box,draw=omThm,fill=omThm!8] (iss) at (4.6,0) {émetteur};
\node[box,draw=omThm,fill=omThm!8] (dl) at (9.2,0) {dealer};
\node[box,draw=omProp,fill=omProp!8] (vb) at (13.2,0) {acheteurs\\de volatilité};
\draw[arr] (inv) -- node[above=5pt,font=\footnotesize,align=center] {achète un\\callable à un\\coupon plus élevé} (iss);
\draw[arr] (iss) -- node[above=5pt,font=\footnotesize,align=center] {annulable\\swap~: vend\\le call} (dl);
\draw[arr] (dl) -- node[above=5pt,font=\footnotesize,align=center] {vend\\des swaptions} (vb);
\node[font=\footnotesize,text=omThm,align=center] at (4.6,-1.2) {paie un taux variable};
\node[font=\footnotesize,text=omThm,align=center] at (9.2,-1.25) {longue d'une swaption receveuse,\\longue de vega jusqu'à ce que le risque soit transféré};
\end{tikzpicture}

\omcaption{Comment l'option incluse dans une \omterm{def:m2:the-rates-options-market:swaption}{obligation remboursable par anticipation} (callable) atteint le marché de la volatilité.
Les investisseurs sont rémunérés pour le call de l'émetteur dans le coupon~; l'émetteur vend la
même option à un dealer au sein d'un swap annulable et paie un taux variable~; le dealer
vend des \omterm{def:m2:the-rates-options-market:swaption}{swaptions} aux investisseurs qui veulent de la volatilité. Les flèches indiquent qui vend
l'option à qui. Schéma.}\label{fig:m2:the-rates-options-market:callable}
\end{omfigure}

Les hypothèques fonctionnent dans l'autre sens. Le détenteur d'un titre adossé à des hypothèques est vendeur
de l'option de remboursement anticipé des emprunteurs (du \cref{ch:m2:mortgages-and-agencies}), et
certains se couvrent contre cette convexité négative en achetant des \omterm{def:m2:the-rates-options-market:swaption}{swaptions}. Les émetteurs
de dettes remboursables par anticipation sont des vendeurs naturels de volatilité, les investisseurs
hypothécaires des acheteurs naturels, et les dealers se placent entre les deux.

\section{Volatilité normale}

\begin{definition}[Modèle de Bachelier, volatilité normale]\label{def:m2:the-rates-options-market:normal}
Le \emph{modèle de Bachelier}\index{modèle de Bachelier}, publié par Louis Bachelier
en 1900, suppose que le taux forward suit $dF_t = \sigma_N\,dW_t$~: il
varie de montants distribués normalement, indépendamment de son niveau, et peut devenir
négatif. Son paramètre $\sigma_N$, exprimé en points de base par an, est la \emph{volatilité
normale}\index{volatilité normale}~; la volatilité normale implicite d'une option
est le $\sigma_N$ auquel le modèle reproduit son prix.
\end{definition}

\begin{proposition}[Le prix de Bachelier]\label{prop:m2:the-rates-options-market:bachelier}
Dans le \omterm{def:m2:the-rates-options-market:normal}{modèle de Bachelier}, un payeur de strike $K$ et d'échéance $T$, sur un taux forward
$F$ d'annuité $A$, vaut
\[
A\left[(F-K)\,\Phi(d) + \sigma_N\sqrt{T}\,\varphi(d)\right], \qquad d = \frac{F-K}{\sigma_N\sqrt T},
\]
et un receveur vaut le payeur diminué de $A(F-K)$. À la monnaie, chacun vaut
$A\sigma_N\sqrt{T}/\sqrt{2\pi}$, de sorte qu'un straddle coûte
$A\sigma_N\sqrt{T}\sqrt{2/\pi}$.
\end{proposition}

\begin{proof}
$F_T$ suit une loi normale de moyenne $F$ et d'écart type $s = \sigma_N\sqrt{T}$.
En posant $F_T = F + sZ$, $\mathbb{E}[(F_T-K)^+] = \mathbb{E}[(F-K+sZ)\mathbf 1\{Z>-d\}]
= (F-K)\Phi(d) + s\varphi(d)$, en utilisant $\mathbb{E}[Z\mathbf 1\{Z>-d\}] = \varphi(d)$.
La parité call-put donne le receveur, et $d = 0$ donne la valeur à la monnaie.
\end{proof}

Un trader convertit la \omterm{def:m2:the-rates-options-market:normal}{volatilité normale} annuelle en volatilité journalière pour la comparer
à ce que font réellement les taux~: 95 points de base par an correspondent à
$95/\sqrt{252} = 5.98$ points de base par jour ouvré. Si le taux à dix ans bouge
en moyenne de plus de six points de base par jour sur la durée de vie de l'option, un
straddle couvert en delta acheté à 95 rapporte~; moins, et il perd. Coter en
\omterm{def:m2:the-rates-options-market:normal}{volatilité normale} rend cette comparaison immédiate, et cela reste la convention usuelle
pour les options sur taux, maintenant que les taux sont redevenus positifs.

La cotation lognormale, dite de Black, reste utilisée et doit être convertie. À la monnaie,
$\sigma_N \approx \sigma_B F$~; ailleurs, les sourires des modèles diffèrent. Près de zéro,
la conversion explose (la \cref{fig:m2:the-rates-options-market:black})~: 90
points de base de \omterm{def:m2:the-rates-options-market:normal}{volatilité normale} correspondent à une volatilité de Black de 22.5~\% pour un
taux forward à 4~\%, 93~\% à 1~\% et 215~\% à 0.5~\%, et en dessous d'un taux forward d'environ
0.36~\% aucune volatilité de Black ne reproduit le prix, car un call de Black ne peut jamais
valoir plus que le forward. Lorsque les taux en euros et en yens ont passé des années
proches de zéro ou en dessous (la \cref{fig:m2:the-rates-options-market:negative}), les cotations
de Black ont cessé de porter de l'information, et la cotation en \omterm{def:m2:the-rates-options-market:normal}{volatilité normale} s'est imposée.

\begin{omfigure}
\begin{tikzpicture}
\begin{axis}[width=11cm,height=4.8cm,ymode=log,
  xlabel={taux forward (\%)},ylabel={volatilité de Black (\%)},
  tick label style={font=\small},label style={font=\small},
  xmin=0,xmax=6.2,ymin=10,ymax=500,grid=major,grid style={gray!25},
  log ticks with fixed point,
  table/col sep=comma]
\addplot[omThm,very thick,mark=*,mark size=1.2pt] table[x=forward,y=black]{figdata/markets-2/13-the-rates-options-market/blackequiv.csv};
\draw[omProp,thick,dashed] (axis cs:0.359,10) -- (axis cs:0.359,500);
\node[omProp,font=\footnotesize,anchor=west,align=left] (nb) at (axis cs:0.9,19) {aucune volatilité de Black\\en dessous de 0.36~\%};
\draw[omProp,thick,-{Stealth}] (nb.west) -- (axis cs:0.40,19);
\end{axis}
\end{tikzpicture}

\omcaption{La volatilité de Black qui donne le même prix à un an à la monnaie
qu'une \omterm{def:m2:the-rates-options-market:normal}{volatilité normale} de 90 points de base, par taux forward, sur une échelle
logarithmique. La même option paraît calme à 5~\% et agitée à 0.5~\%~; sous la ligne
pointillée, aucune volatilité de Black n'existe. Données~: le tutoriel du
chapitre.}\label{fig:m2:the-rates-options-market:black}
\end{omfigure}

\begin{omfigure}
\begin{tikzpicture}
\begin{axis}[width=11cm,height=4.6cm,
  ylabel={rendement à dix ans (\%)},
  tick label style={font=\small},label style={font=\small},
  xmin=2012,xmax=2023,ymin=-0.8,ymax=2.2,grid=major,grid style={gray!25},
  xtick={2012,2014,2016,2018,2020,2022},xticklabel style={/pgf/number format/1000 sep={}},
  legend columns=2,legend style={font=\footnotesize,at={(0.5,-0.2)},anchor=north,draw=none,fill=none,/tikz/every even column/.append style={column sep=8pt}},
  table/col sep=comma]
\addplot[omDef,very thick] table[x=t,y=DE]{figdata/markets-2/13-the-rates-options-market/negyields.csv};
\addplot[omThm,thick,dashed] table[x=t,y=JP]{figdata/markets-2/13-the-rates-options-market/negyields.csv};
\draw[gray] (axis cs:2012,0) -- (axis cs:2023,0);
\legend{Allemagne,Japon}
\end{axis}
\end{tikzpicture}

\omcaption{Rendements des obligations d'État à dix ans allemandes et japonaises, moyennes mensuelles,
de 2012 à 2022. Le rendement allemand a été négatif pendant 38 mois entre juin 2016 et
janvier 2022, le japonais pendant 24 mois entre février 2016 et avril 2020. Données~:
OCDE via les séries FRED IRLTLT01DEM156N et
IRLTLT01JPM156N.}\label{fig:m2:the-rates-options-market:negative}
\end{omfigure}

\begin{example}[Un straddle]\label{ex:m2:the-rates-options-market:straddle}
Sur une courbe plate à 4~\%, l'annuité forward 1a$\times$10a vaut 7.80. Avec une \omterm{def:m2:the-rates-options-market:normal}{volatilité normale}
de 95 points de base, un straddle sur 100~millions de dollars américains coûte
$7.80 \times 0.0095 \times \sqrt{2/\pi} = 5.91\%$ du notionnel, soit 5.91~millions de dollars américains.
Divisé par l'annuité, cela fait 75.8 points de base~: le taux à dix ans doit finir
à plus de 76 points de base environ de 4~\% pour que le straddle conservé jusqu'à l'échéance
se rembourse.
\end{example}

\section{Le cube de volatilité et ceux qui le négocient}

\begin{definition}[Cube de volatilité]\label{def:m2:the-rates-options-market:cube}
Le \emph{cube de volatilité}\index{cube de volatilité} d'une devise est l'ensemble des
volatilités implicites de ses \omterm{def:m2:the-rates-options-market:swaption}{swaptions}, par échéance, par maturité du swap sous-jacent et
par prix d'exercice. Sa face à la monnaie, échéance par maturité, est la matrice des \omterm{def:m2:the-rates-options-market:swaption}{swaptions}~; la
dimension des prix d'exercice est le sourire.
\end{definition}

Un desk d'options de taux cote et gère le risque de l'ensemble du cube. Sa matrice à la monnaie
possède une structure par terme dans chaque direction
(\cref{fig:m2:the-rates-options-market:cube})~: les échéances courtes sur les maturités courtes
évoluent avec les prochaines décisions des banques centrales, les échéances longues sur les maturités longues
avec les flux lents des émissions remboursables par anticipation et des couvertures hypothécaires. Le sourire, que
un modèle comme SABR (livre~6) interpole, price les queues~; pour les taux
il est en général plus élevé des deux côtés de la monnaie, et asymétrique.

\begin{omfigure}
\begin{tikzpicture}
\begin{axis}[name=left,width=6.4cm,height=4.8cm,xmode=log,
  xlabel={échéance (années)},ylabel={volatilité normale (pb par an)},
  tick label style={font=\small},label style={font=\small},
  xmin=0.2,xmax=12,ymin=70,ymax=120,grid=major,grid style={gray!25},
  xtick={0.25,1,5,10},xticklabels={3m,1a,5a,10a},
  legend columns=3,legend style={font=\footnotesize,at={(1.08,-0.3)},anchor=north,draw=none,fill=none,/tikz/every even column/.append style={column sep=8pt}},
  table/col sep=comma]
\addplot[omDef,very thick,mark=*,mark size=1.2pt] table[x=expiry,y=t2]{figdata/markets-2/13-the-rates-options-market/cube.csv};
\addplot[omThm,very thick,mark=*,mark size=1.2pt] table[x=expiry,y=t10]{figdata/markets-2/13-the-rates-options-market/cube.csv};
\addplot[omProp,very thick,mark=*,mark size=1.2pt] table[x=expiry,y=t30]{figdata/markets-2/13-the-rates-options-market/cube.csv};
\legend{maturité 2~ans,maturité 10~ans,maturité 30~ans}
\end{axis}
\begin{axis}[at={(left.east)},anchor=west,xshift=1.4cm,width=6.4cm,height=4.8cm,
  xlabel={prix d'exercice moins forward (pb)},
  tick label style={font=\small},label style={font=\small},
  xmin=-210,xmax=210,ymin=90,ymax=112,grid=major,grid style={gray!25},
  table/col sep=comma]
\addplot[omThm,very thick,mark=*,mark size=1.2pt] table[x=offset,y=vol]{figdata/markets-2/13-the-rates-options-market/smile.csv};
\end{axis}
\end{tikzpicture}

\omcaption{Deux faces d'un \omterm{def:m2:the-rates-options-market:cube}{cube de volatilité} en dollars, à titre illustratif~: \omterm{def:m2:the-rates-options-market:normal}{volatilité normale} à la
monnaie par échéance pour trois maturités (à gauche), et le sourire de la
\omterm{def:m2:the-rates-options-market:swaption}{swaption} 1a$\times$10a par prix d'exercice (à droite). Ce sont les formes, et non les chiffres, qui
comptent. Données~: le tutoriel du chapitre.}\label{fig:m2:the-rates-options-market:cube}
\end{omfigure}

Les flux qui façonnent le cube peuvent être importants et unidirectionnels. À Taïwan, une loi de
mai 2014 a permis aux assureurs-vie de détenir des obligations libellées en devises étrangères émises localement, hors de
leur plafond sur les actifs à l'étranger~; chaque obligation en dollars vendue là-bas dans les mois suivants
était remboursable par anticipation par son émetteur, avec des maturités de vingt ou trente ans,
et trois sur quatre ne versaient aucun coupon. Chacune de ces émissions, swappée par son émetteur, laisse
un dealer long d'une \omterm{def:m2:the-rates-options-market:swaption}{swaption receveuse} de longue échéance qu'il doit revendre~: une offre régulière
de volatilité longue, issue des bilans
des assureurs d'un seul pays.

\begin{dated}{2026-09}{La taille du marché}\label{dat:m2:the-rates-options-market:size}
L'encours notionnel des dérivés de gré à gré s'élevait à 846~billions de dollars américains à la fin
de juin 2025, soit 16~\% de plus qu'un an plus tôt, dont 79~\% sur les dérivés de taux~; leur valeur de
marché brute s'élevait à 21.8~billions de dollars américains (statistiques de la BRI
publiées le 8 décembre 2025). Les \omterm{def:m2:the-rates-options-market:swaption}{swaptions}, \omterm{def:m2:the-rates-options-market:cap}{caps} et \omterm{def:m2:the-rates-options-market:cap}{floors} font partie du total des taux~;
la BRI détaille les options dans ses tableaux détaillés.
\end{dated}

\section{Tutoriel~: cotations normales et lognormales}\label{tut:m2:the-rates-options-market:quotes}

\begin{tutorial}
\textbf{Objectif.} Pricer des \omterm{def:m2:the-rates-options-market:cap}{caps} et des \omterm{def:m2:the-rates-options-market:swaption}{swaptions} dans les \omterm{def:m2:the-rates-options-market:normal}{modèles de Bachelier} et de Black,
convertir des cotations de l'un à l'autre, trouver où la conversion échoue, et pricer un
straddle et l'option d'une \omterm{def:m2:the-rates-options-market:swaption}{obligation remboursable par anticipation} à partir d'une tranche du cube. \textbf{État final~:}
\cref{fig:m2:the-rates-options-market:payoff,fig:m2:the-rates-options-market:black},
\cref{ex:m2:the-rates-options-market:cap,ex:m2:the-rates-options-market:straddle}
et les chiffres du problème du week-end.
\begin{tutsteps}
\item \textbf{Les deux pricers et le vega} de
la \cref{prop:m2:the-rates-options-market:bachelier}.
\omcode{code/firm/normalvol/firm_normalvol.py}{20}{43}{Pricers de Bachelier et de Black, vega normal.}\label{lst:m2:the-rates-options-market:pricers}
\item \textbf{Volatilité implicite et conversion}, par dichotomie, en refusant un
prix qu'aucune volatilité ne reproduit.
\omcode{code/firm/normalvol/firm_normalvol.py}{46}{68}{Volatilité implicite et conversion des cotations.}\label{lst:m2:the-rates-options-market:implied}
\item \textbf{Exécuter} \texttt{swaption\_demo.straddle\_1y10y()},
\texttt{swaption\_demo.cap()}, \texttt{swaption\_demo.callable\_swap()} et
\texttt{fig\_swaptions.py}.
\end{tutsteps}
\textbf{Pour aller plus loin.} Convertir le sourire de
la \cref{fig:m2:the-rates-options-market:cube} en volatilités de Black à un forward de 4~\%
et à un forward de 1~\%, et comparer les formes~; puis décaler la courbe vers le bas de 3.5
points et recommencer.
\end{tutorial}

\section{Construction~: le pricer de volatilité normale}\label{bld:m2:the-rates-options-market:normalvol}

\begin{build}
\bfield{Rôle} La mini-société cote des \omterm{def:m2:the-rates-options-market:cap}{caps}, des \omterm{def:m2:the-rates-options-market:cap}{floors} et des \omterm{def:m2:the-rates-options-market:swaption}{swaptions} en
\omterm{def:m2:the-rates-options-market:normal}{volatilité normale}, reçoit des cotations de Black de certaines contreparties et doit les
convertir, et couvre le vega issu de transactions structurées.
\bfield{Interface} \texttt{bachelier(f, k, t, vol, annuity, payer)}~;
\texttt{black(\ldots)}~; \texttt{normal\_vega}~; \texttt{implied(price, f, k, t,
annuity, payer, model)}~; \texttt{normal\_to\_black}~;
\texttt{black\_to\_normal}~; \texttt{atm\_straddle(t, vol, annuity)}~;
\texttt{bp\_per\_day(vol)}.
\bfield{Règles} Taux et volatilités en décimales~; Black refuse un forward ou un prix d'exercice
non positif~; volatilité implicite par dichotomie entre des bornes explicites,
avec une erreur levée lorsque le prix est hors de l'intervalle du modèle.
\bfield{Tests d'acceptation} \texttt{code/firm/normalvol/tests/}~: parité call-put dans
les deux modèles, y compris avec des taux négatifs~; la forme fermée à la monnaie~; le vega
comparé à une différence finie~; les allers-retours de volatilité implicite et de conversion~;
le refus à un forward de 0.25~\%.
\bfield{Pour aller plus loin} Black décalé~; SABR sous sa forme normale (livre~6)~;
\omterm{def:m2:the-rates-options-market:swaption}{swaptions} bermudéennes par un arbre (livre~6)~; un interpolateur de cube par
échéance et maturité.
\end{build}

\begin{omsources}
\item L.~Bachelier, ``Théorie de la spéculation'', \emph{Annales
scientifiques de l'École Normale Supérieure}, 1900.
\item Banque centrale européenne, taux directeurs de la BCE~; taux d'intérêt à long terme
de l'OCDE via FRED.
\item F.~Le Floc'h, ``Explicit rational formulae for Bachelier (normal) implied
volatility'', arXiv 2605.18343, 2026.
\item Bloomberg, ``Goldman leads Wall Street tapping Taiwan cash pile'',
10~septembre 2014 (via Insurance Journal).
\item Banque des règlements internationaux, OTC derivatives statistics at
end-June 2025.
\end{omsources}

\section{Exercices}

\begin{exercise}[$\star$]\label{exo:m2:the-rates-options-market:1}
Exprimer une \omterm{def:m2:the-rates-options-market:normal}{volatilité normale} de 95 points de base par an en points de base par
jour ouvré, et une volatilité journalière de 7 points de base sous forme annuelle.
\end{exercise}

\begin{exercise}[$\star$]\label{exo:m2:the-rates-options-market:2}
Une entreprise a un emprunt à taux variable. Quelle option achète-t-elle pour se protéger, une \omterm{def:m2:the-rates-options-market:swaption}{swaption payeuse}
ou une \omterm{def:m2:the-rates-options-market:swaption}{swaption receveuse}, un \omterm{def:m2:the-rates-options-market:cap}{cap} ou un \omterm{def:m2:the-rates-options-market:cap}{floor}~? Et un fonds de pension qui craint une baisse
des taux~?
\end{exercise}

\begin{exercise}[$\star$]\label{exo:m2:the-rates-options-market:3}
Pourquoi le modèle de Black ne peut-il pas évaluer une option sur un taux de $-0.2\%$, et
pourquoi le marché a-t-il choisi la \omterm{def:m2:the-rates-options-market:normal}{volatilité normale} plutôt qu'une autre correction~?
\end{exercise}

\begin{exercise}[$\star\star$]\label{exo:m2:the-rates-options-market:4}
Donner le prix du payeur à la monnaie 1an$\times$10ans sur 100 millions de dollars à
95 points de base, puis le mouvement d'équilibre du straddle.
\end{exercise}

\begin{exercise}[$\star\star$]\label{exo:m2:the-rates-options-market:5}
Donner les volatilités de Black équivalentes à une \omterm{def:m2:the-rates-options-market:normal}{volatilité normale} de 90 points de base,
à un an à la monnaie, pour des taux forward de 6~\%, 2~\% et 0.4~\%.
\end{exercise}

\begin{exercise}[$\star\star$]\label{exo:m2:the-rates-options-market:6}
Dans l'\cref{ex:m2:the-rates-options-market:cap}, pourquoi le \omterm{def:m2:the-rates-options-market:cap}{floor} à 3.5~\% vaut-il
exactement le \omterm{def:m2:the-rates-options-market:cap}{cap} à 4.5~\%~? Serait-il encore ainsi dans le modèle de Black~?
\end{exercise}

\begin{exercise}[$\star\star\star$]\label{exo:m2:the-rates-options-market:7}
\emph{Programmation.} Trouver le forward en dessous duquel aucune volatilité de Black ne reproduit
le prix à la monnaie à un an pour une \omterm{def:m2:the-rates-options-market:normal}{volatilité normale} de 90 points de base, puis le vérifier
par rapport à $\sigma_N\sqrt T/\sqrt{2\pi}$.
\end{exercise}

\begin{exercise}[$\star\star\star$]\label{exo:m2:the-rates-options-market:8}
\emph{Repérer l'erreur.} «~La volatilité des \omterm{def:m2:the-rates-options-market:swaption}{swaptions} est de 20~\% dans une devise et de 60~\% dans
une autre~: la seconde est trois fois plus risquée.~» La corriger.
\end{exercise}

\section{Problème~: L'émetteur remboursable}

\begin{problem}[{Problème du week-end --- la vega que prend une banque sur une obligation remboursable par anticipation}]\label{pb:m2:the-rates-options-market:1}
Une banque émet 500 millions de dollars d'obligations à dix ans portant un coupon annuel de 4.5~\%,
remboursables au pair une fois, dans trois ans. Elle les échange auprès d'un dealer
contre un taux variable~: le dealer verse un taux fixe de 4.5~\% et reçoit le taux variable pendant dix ans, et
peut annuler le swap dans trois ans, lorsque la banque rembourserait l'obligation. La
courbe est plate à 4~\%~; la \omterm{def:m2:the-rates-options-market:normal}{volatilité normale} 3ans$\times$7ans est de 100 points de base.

\medskip\noindent\textbf{Partie I --- L'option.}
\begin{enumerate}
\item Quelle option la banque détient-elle à travers son obligation, et pourquoi~?
\item Quelle option le dealer détient-il à travers son swap~?
\item Donner le taux de swap forward et l'annuité forward du swap 3ans$\times$7ans.
\item L'option est-elle dans la monnaie ou hors de la monnaie, et de combien~?
\item Donner sa valeur sur 500 millions de dollars, en dollars et en points de base du notionnel.
\end{enumerate}

\medskip\noindent\textbf{Partie II --- Le risque du dealer.}
\begin{enumerate}[resume]
\item Donner sa vega par point de base de \omterm{def:m2:the-rates-options-market:normal}{volatilité normale}.
\item Donner son delta par point de base du taux de swap forward, et son signe.
\item De combien le dealer perd-il si la volatilité baisse de 2 points de base~?
\item Donner la vega de 100 millions de dollars de straddles 3ans$\times$7ans à la monnaie.
\item Quel notionnel de straddles le dealer doit-il vendre pour couvrir la vega~?
\end{enumerate}

\medskip\noindent\textbf{Partie III --- La banque.}
\begin{enumerate}[resume]
\item Comment la banque est-elle rémunérée pour l'option qu'elle vend~?
\item Exprimer la valeur de l'option comme un spread courant sur les dix ans.
\item Quand la banque rembourse-t-elle l'obligation, et que devient le swap~?
\item Pourquoi la banque préfère-t-elle cela à une obligation simple et à un swap simple~?
\item Que devient la couverture du dealer si les taux baissent à 3.5~\% le lendemain~?
\end{enumerate}

\medskip\noindent\textbf{Partie IV --- Le marché.}
\begin{enumerate}[resume]
\item Que se passe-t-il pour la volatilité des \omterm{def:m2:the-rates-options-market:swaption}{swaptions} de longue échéance lorsque de nombreux émetteurs font de même~?
\item Qui achète la volatilité que les dealers vendent~?
\item Pourquoi une véritable obligation remboursable chaque année est-elle plus difficile à couvrir~?
\item Énoncer le \emph{résultat nommé}~: la vega que le dealer acquiert, et les
straddles qu'il vend.
\item En une phrase~: d'où vient l'offre de volatilité de longue échéance~?
\end{enumerate}
\end{problem}

\section{Questions d'entretien}

\begin{interviewq}[$\star$ \iqroles{trader, chercheur}]\label{iq:m2:the-rates-options-market:1}
Quelle est la différence entre \omterm{def:m2:the-rates-options-market:normal}{volatilité normale} et volatilité lognormale, et quand
cela compte-t-il~?
\end{interviewq}

\begin{interviewq}[$\star$ \iqroles{trader, banque}]\label{iq:m2:the-rates-options-market:2}
Expliquez ce qu'est une \omterm{def:m2:the-rates-options-market:swaption}{swaption payeuse} et une \omterm{def:m2:the-rates-options-market:swaption}{swaption receveuse}, et qui achète chacune.
\end{interviewq}

\begin{interviewq}[$\star\star$ \iqroles{chercheur}]\label{iq:m2:the-rates-options-market:3}
Dérivez le prix d'une option à la monnaie dans le \omterm{def:m2:the-rates-options-market:normal}{modèle de Bachelier}.
\end{interviewq}

\begin{interviewq}[$\star\star$ \iqroles{trader}]\label{iq:m2:the-rates-options-market:4}
Vous êtes long d'un straddle 1an$\times$10ans couvert en delta à 95 points de base. Comment
savez-vous, jour après jour, si vous gagnez de l'argent~?
\end{interviewq}

\begin{interviewq}[$\star\star$ \iqroles{banque, trader}]\label{iq:m2:the-rates-options-market:5}
Pourquoi l'émission d'obligations remboursables fait-elle baisser la volatilité des \omterm{def:m2:the-rates-options-market:swaption}{swaptions} de longue échéance, et
la couverture des prêts hypothécaires la fait-elle monter~?
\end{interviewq}

\begin{interviewq}[$\star\star\star$ \iqroles{développeur, chercheur}]\label{iq:m2:the-rates-options-market:6}
La \omterm{def:m2:the-rates-options-market:normal}{volatilité normale} implicite doit être calculée pour un million d'options par seconde.
Comment le feriez-vous, et comment le testeriez-vous~?
\end{interviewq}
